% Contributions:
% 1. We introduce the concept of flatness into domain generalization
% 2. We propose SWAD to capture flatter minima
% 3. SWAD achieves state-of-the-art performances with large margin
% 4. We show even better results by minimizing both robust risk and domain divergence

% \jb{ToDo: Update conclusion also before finalize paper.}

In this paper, we theoretically and empirically demonstrate that domain generalization (DG) is achievable by seeking flat minima.
We propose SWAD that captures flatter minima than the vanilla SWA does.
The extensive experiments on five DG benchmarks show superior performances of SWAD compared with existing DG methods.
In addition, combinations of SWAD and existing DG methods even show better performances than the vanilla SWAD.
We theoretically and empirically observe that seeking flat minima can achieve better generalizability to both in-domain and out-of-domain, while strong in-domain generalization methods without consideration of flatness, \eg, Mixup or CutMix, cannot guarantee to achieve out-of-domain generalizability in both theory and practice.
This study first brings the concept of flatness into DG tasks, and shows strong empirical performances not only in DG but also in ImageNet benchmarks.
% As our theoretical observation supports, seeking flat minima can achieve better DG, while strong in-domain generalization without consideration of flatness, \eg, Mixup or CutMix, does not have any guarantee.
% Comparisons with conventional generalization methods, such as Mixup or CutMix, supports that the significant improvement by SWAD is from seeking flat minima, not from better in-domain generalization performances.
% It also supports our theoretical result that the DG gap is bounded to robust risk and divergence between training and target domains.
% It should be noted that this study is very first work to introduce the concept of flatness into the DG field.
We hope that this study promotes a new research direction of seeking flat minima for domain generalization and other robustness tasks.
% We believe this study is a foundation work in a new research direction of seeking flat minima for domain generalization.
% We believe this study open a new direction, seeking flat minima, in DG field, and hope that this study will serve as cornerstone 

% In addition, we show that combinations of SWAD and other DG methods can achieve state-of-the-art performances on the DG benchmarks.
